% =============================================================================
%  Relatório Técnico — Projecto Nº 1
%  Engenharia de Reservatórios I — 2025/2026
%  ISPTEC — Instituto Superior Politécnico de Tecnologias e Ciências
%  Docente: Geraldo Ramos, BSc, MSc, PhD
% =============================================================================
\documentclass[12pt,a4paper]{article}

% ── Codificação e língua ──────────────────────────────────────────────────────
\usepackage[utf8]{inputenc}
\usepackage[T1]{fontenc}
\usepackage[portuguese]{babel}

% ── Geometria ─────────────────────────────────────────────────────────────────
\usepackage{geometry}
\geometry{a4paper, left=3cm, right=2cm, top=3cm, bottom=2cm}

% ── Tipografia ────────────────────────────────────────────────────────────────
\usepackage{times}
\usepackage{microtype}
\usepackage{setspace}
\onehalfspacing
\setlength{\parindent}{1.25cm}
\setlength{\parskip}{0pt}

% ── Matemática ────────────────────────────────────────────────────────────────
\usepackage{amsmath}
\usepackage{amssymb}
\usepackage[version=4]{mhchem}
\usepackage{siunitx}
\sisetup{locale=FR, output-decimal-marker={,}}

% ── Figuras e tabelas ─────────────────────────────────────────────────────────
\usepackage{graphicx}
\usepackage{booktabs}
\usepackage{longtable}
\usepackage{array}
\usepackage{multirow}
\usepackage{caption}
\usepackage{float}
\captionsetup{font=small, labelfont=bf, justification=centering, skip=4pt}

% ── Código-fonte ──────────────────────────────────────────────────────────────
\usepackage{listings}
\usepackage{xcolor}
\definecolor{codebg}{RGB}{245,245,245}
\definecolor{codekey}{RGB}{0,0,180}
\definecolor{codecom}{RGB}{60,128,49}
\definecolor{codestr}{RGB}{163,21,21}
\lstset{
  language=Python,
  backgroundcolor=\color{codebg},
  basicstyle=\ttfamily\small,
  keywordstyle=\color{codekey}\bfseries,
  commentstyle=\color{codecom}\itshape,
  stringstyle=\color{codestr},
  numbers=left, numberstyle=\tiny\color{gray}, numbersep=8pt,
  breaklines=true, breakatwhitespace=false,
  tabsize=4, showspaces=false, showstringspaces=false,
  frame=single, framesep=3pt, rulecolor=\color{gray!40},
  captionpos=b
}

% ── Hiperligações ─────────────────────────────────────────────────────────────
\usepackage{hyperref}
\hypersetup{
  colorlinks=true, linkcolor=blue!60!black,
  citecolor=blue!60!black, urlcolor=blue!60!black,
  pdftitle={Factor de Compressibilidade do Gás Natural - Projecto 1},
  pdfauthor={ISPTEC - Engenharia de Reservatórios I}
}

% ── Cabeçalho / rodapé ────────────────────────────────────────────────────────
\usepackage{fancyhdr}
\pagestyle{fancy}
\fancyhf{}
\fancyhead[L]{\small\textit{Engenharia de Reservatórios I — Projecto Nº 1}}
\fancyhead[R]{\small\textit{ISPTEC — 2025/2026}}
\fancyfoot[C]{\small\thepage}
\renewcommand{\headrulewidth}{0.4pt}

% ── Utilitários ───────────────────────────────────────────────────────────────
\usepackage{enumitem}
\usepackage{lastpage}
\usepackage{indentfirst}

% =============================================================================
%  CAPA
% =============================================================================
\begin{document}
\thispagestyle{empty}

\begin{center}
  {\Large\bfseries INSTITUTO SUPERIOR POLITÉCNICO DE\\[4pt]
  TECNOLOGIAS E CIÊNCIAS}\\[6pt]
  {\large ISPTEC}\\[1.5cm]

  {\large\bfseries ENGENHARIA DE RESERVATÓRIOS I}\\[4pt]
  {\normalsize 3.º Ano — 6.º Semestre — 2025/2026}\\[2cm]

  \rule{\linewidth}{1pt}\\[0.4cm]
  {\LARGE\bfseries Cálculo do Factor de Compressibilidade\\[6pt]
  do Gás Natural (Factor~\textit{Z})}\\[0.3cm]
  {\large Projecto Nº~1}\\[0.2cm]
  \rule{\linewidth}{1pt}\\[2cm]

  {\normalsize
  \begin{tabular}{rl}
    \textbf{Docente:}  & Geraldo Ramos, BSc, MSc, PhD \\[4pt]
    \textbf{Unidade:}  & Engenharia de Reservatórios I \\[4pt]
    \textbf{Data:}     & Maio de 2026 \\[4pt]
    \textbf{Prazo:}    & 14 de Maio de 2026 \\
  \end{tabular}}

  \vfill
  {\small Luanda, Angola — 2026}
\end{center}

\newpage
\pagenumbering{roman}
\setcounter{page}{1}

% =============================================================================
%  RESUMO
% =============================================================================
\section*{Resumo}
\addcontentsline{toc}{section}{Resumo}

Este relatório descreve o desenvolvimento de um programa computacional em Python
para o cálculo do factor de compressibilidade do gás natural~($Z$), no âmbito
do Projecto~Nº~1 da unidade curricular de Engenharia de Reservatórios~I
(2025/2026) do ISPTEC.

O programa implementa três metodologias de cálculo do factor~$Z$: gás ideal
($Z=1$), correlação de Hall-Yarborough com resolução iterativa pelo método de
Newton-Raphson, e correlação de Dranchuk \& Abou-Kassem com iteração de ponto
fixo. Adicionalmente, são implementadas correlações de propriedades
pseudo-críticas (Standing e Sutton) com correcções para contaminantes
(Wichert-Aziz e Carr-Kobayashi-Burrows), bem como o cálculo da viscosidade do
gás, do factor volume de formação ($B_g$) e do factor de expansão ($E_g$).

O programa foi validado com dados de referência obtidos nas notas de aula~\cite{Ramos2020}, com
boa concordância entre os métodos Hall-Yarborough e Dranchuk-Abou-Kassem
(divergência inferior a 0{,}3\,\%). A interface gráfica moderna, desenvolvida
com \textit{tkinter}/\textit{ttk} e \textit{matplotlib}, permite a introdução
interactiva de dados, visualização de resultados em tabela e geração de gráficos.

\medskip
\noindent\textbf{Palavras-chave:} factor~$Z$, gás natural, Hall-Yarborough,
Dranchuk-Abou-Kassem, Newton-Raphson, propriedades pseudo-críticas, Python.

\newpage

% =============================================================================
%  ÍNDICE
% =============================================================================
\tableofcontents
\newpage
\pagenumbering{arabic}
\setcounter{page}{1}

% =============================================================================
%  1. INTRODUÇÃO
% =============================================================================
\section{Introdução}

O comportamento dos gases reais em condições de reservatório difere
significativamente do previsto pela lei dos gases ideais, especialmente a
elevadas pressões e temperaturas. Esta diferença é quantificada pelo
\emph{factor de compressibilidade}~$Z$ (também denominado factor de desvio),
definido como:

\begin{equation}
  Z = \frac{PV}{nRT}
  \label{eq:z_def}
\end{equation}

\noindent onde $P$ é a pressão, $V$ o volume, $n$ o número de moles, $R$ a
constante universal dos gases e $T$ a temperatura absoluta. Para um gás ideal,
$Z=1$ em qualquer condição. Para gases reais, $Z$ pode ser inferior ou superior
a~1 dependendo das condições de pressão e temperatura.

O factor~$Z$ é fundamental na engenharia de reservatórios, pois afecta
directamente:
\begin{itemize}[noitemsep]
  \item o cálculo das reservas de gás (\emph{gas in place});
  \item o factor volume de formação do gás ($B_g$);
  \item a estimativa da produção e comportamento do reservatório.
\end{itemize}

\subsection{Objectivos do Projecto}

Este projecto tem como objectivos pedagógicos:
\begin{enumerate}[noitemsep]
  \item Desenvolver competências em modelação computacional aplicada a reservatórios;
  \item Implementar as principais correlações empíricas do factor~$Z$;
  \item Comparar os resultados dos diferentes métodos;
  \item Estender o programa ao cálculo de propriedades derivadas ($\mu_g$, $B_g$, $E_g$).
\end{enumerate}

% =============================================================================
%  2. FUNDAMENTAÇÃO TEÓRICA
% =============================================================================
\section{Fundamentação Teórica}

\subsection{Propriedades Pseudo-Críticas}

As correlações do factor~$Z$ utilizam a \emph{pressão pseudo-reduzida}
($P_{pr}$) e a \emph{temperatura pseudo-reduzida}~($T_{pr}$), definidas por:

\begin{equation}
  P_{pr} = \frac{P}{P_{pc}}, \qquad T_{pr} = \frac{T}{T_{pc}}
  \label{eq:reduzidas}
\end{equation}

\noindent onde $P_{pc}$ e $T_{pc}$ são as propriedades pseudo-críticas da
mistura, estimadas pelas correlações apresentadas nas subsecções seguintes.

\subsubsection{Correlação de Standing (1977) — Gás Natural Seco}

Standing~\cite{Standing1977} propôs as seguintes correlações para gás natural seco:
\begin{align}
  P_{pc} &= 677{,}0 + 15{,}0\,\gamma_g - 37{,}5\,\gamma_g^2 \quad [\text{psia}]
  \label{eq:standing_seco_ppc}\\
  T_{pc} &= 168{,}0 + 325{,}0\,\gamma_g - 12{,}5\,\gamma_g^2 \quad [°\text{R}]
  \label{eq:standing_seco_tpc}
\end{align}

\subsubsection{Correlação de Standing (1977) — Gás Natural Húmido/Condensado}

Para gases húmidos e condensados de gás, Standing~\cite{Standing1977} fornece:
\begin{align}
  P_{pc} &= 706{,}0 - 51{,}7\,\gamma_g - 11{,}1\,\gamma_g^2 \quad [\text{psia}]\\
  T_{pc} &= 187{,}0 + 330{,}0\,\gamma_g - 71{,}5\,\gamma_g^2 \quad [°\text{R}]
\end{align}

\subsubsection{Correlação de Sutton (1985)}

A correlação de Sutton é aplicável a gases húmidos e condensados de gás.
As equações correctas, de acordo com o artigo original~\cite{Sutton1985} (SPE-14265), são:
\begin{align}
  P_{pc} &= 756{,}8 - 131{,}07\,\gamma_g - 3{,}6\,\gamma_g^2 \quad [\text{psia}]\\
  T_{pc} &= 169{,}2 + 349{,}5\,\gamma_g - 74{,}0\,\gamma_g^2 \quad [°\text{R}]
\end{align}

\textbf{Nota:} O enunciado do projecto apresenta os rótulos $P_{pc}$ e $T_{pc}$
trocados por erro de digitalização. As equações acima correspondem ao artigo
original e foram verificadas com os dados de referência da figura do enunciado.

\subsection{Correcção de Wichert e Aziz (1972)}

Quando o gás contém \ce{CO2} e/ou \ce{H2S}, Wichert e Aziz~\cite{WichertAziz1972} propuseram as seguintes correcções às propriedades pseudo-críticas:

\begin{align}
  S &= y_{CO_2} + y_{H_2S} \\
  D &= \sqrt{y_{H_2S}} - \left(y_{CO_2}\right)^4 \\
  \varepsilon &= 120\left(S^{0{,}9} - S^{1{,}6}\right) + 15\,D \\
  T_{pc}' &= T_{pc} - \varepsilon \\
  P_{pc}' &= \frac{P_{pc}\,T_{pc}'}{T_{pc} + y_{H_2S}(1-y_{H_2S})\,\varepsilon}
\end{align}

\subsection{Correcção de Carr, Kobayashi e Burrows (1954)}

Como alternativa, Carr, Kobayashi e Burrows~\cite{CarrKB1954} propuseram factores
aditivos directos sobre as propriedades pseudo-críticas:

\begin{align}
  T_{pc}' &= T_{pc} - 80\,y_{CO_2} + 130\,y_{H_2S} - 250\,y_{N_2} \\
  P_{pc}' &= P_{pc} + 440\,y_{CO_2} + 600\,y_{H_2S} - 170\,y_{N_2}
\end{align}

\subsection{Correlação de Hall-Yarborough}

A correlação de Hall-Yarborough~\cite{HallYarborough1974} resolve a equação de
estado de Starling-Carnahan de forma implícita em~$y$ (parâmetro relacionado com
a densidade reduzida):

\begin{equation}
  Z = \frac{0{,}06125\,P_{pr}\,t\,e^{-1{,}2(1-t)^2}}{y}, \qquad t = \frac{1}{T_{pr}}
  \label{eq:Z_HY}
\end{equation}

O valor de~$y$ é determinado iterativamente resolvendo $f(y)=0$:
\begin{equation}
  f(y) = a + b + c + d = 0
\end{equation}
\begin{align}
  a &= -0{,}06125\,P_{pr}\,t\,e^{-1{,}2(1-t)^2}\\
  b &= \frac{(-y+1)y^3 + y^3 + y^2 + y}{(1-y)^3}
     = \frac{\bigl[((-y+1)y+1)y+1\bigr]y}{(1-y)^3}\\
  c &= -\bigl[(4{,}58\,t - 9{,}76)t + 14{,}76\bigr]\,t\,y^2\\
  d &= \bigl[(42{,}4\,t - 242{,}2)t + 90{,}7\bigr]\,t\,y^{2{,}18+2{,}82t}
\end{align}

A derivada~$f'(y)$, necessária ao método de Newton-Raphson, é:
\begin{align}
  f'(y) &= a' + b' + c'\\
  a' &= \frac{\bigl[(((y-4)y+4)y+4)y+1\bigr]}{(1-y)^4}\\
  b' &= -\bigl[(9{,}16\,t - 19{,}52)t + 29{,}52\bigr]\,t\,y\\
  c' &= (2{,}18+2{,}82\,t)\bigl[(42{,}4\,t-242{,}2)t+90{,}7\bigr]\,t\,y^{1{,}18+2{,}82t}
\end{align}

\paragraph{Algoritmo de Newton-Raphson:}
\begin{enumerate}[noitemsep]
  \item Condição inicial: $y_0 = 0{,}001$.
  \item Calcular $f(y)$ e $f'(y)$.
  \item Calcular $y^* = y - f(y)/f'(y)$.
  \item Se $|y^* - y| \geq 10^{-5}$, actualizar $y \leftarrow y^*$ e repetir.
  \item Calcular $Z$ pela equação~\eqref{eq:Z_HY}.
\end{enumerate}

\subsection{Correlação de Dranchuk e Abou-Kassem (1975)}

A correlação de Dranchuk e Abou-Kassem~\cite{DAK1975} baseia-se na equação
de Benedict-Webb-Rubin modificada:

\begin{equation}
  Z = 1 + X_1\,\rho_r + X_2\,\rho_r^2 - X_3\,\rho_r^5
    + X_4\,(1 + A_{11}\,\rho_r^2)\,\frac{\rho_r^2}{T_{pr}^3}\,e^{-A_{11}\,\rho_r^2}
  \label{eq:DAK}
\end{equation}

onde $\rho_r = 0{,}27\,P_{pr}/(Z\,T_{pr})$ é a densidade pseudo-reduzida e:
\begin{align}
  X_1 &= A_1 + \frac{A_2}{T_{pr}} + \frac{A_3}{T_{pr}^3}
         + \frac{A_4}{T_{pr}^4} + \frac{A_5}{T_{pr}^5}\\
  X_2 &= A_6 + \frac{A_7}{T_{pr}} + \frac{A_8}{T_{pr}^2}\\
  X_3 &= A_9\!\left(\frac{A_7}{T_{pr}} + \frac{A_8}{T_{pr}^2}\right)
\end{align}

Os coeficientes ajustados com dados experimentais são:

\begin{table}[H]
\centering
\caption{Coeficientes da correlação de Dranchuk e Abou-Kassem.}
\begin{tabular}{cccccc}
\toprule
$A_1$ & $A_2$ & $A_3$ & $A_4$ & $A_5$ & $A_6$ \\
\midrule
0{,}3265 & $-$1{,}0700 & $-$0{,}5339 & 0{,}01569 & $-$0{,}05165 & 0{,}5475 \\
\midrule
$A_7$ & $A_8$ & $A_9$ & $A_{10}$ & $A_{11}$ & \\
\midrule
$-$0{,}7361 & 0{,}1844 & 0{,}1056 & 0{,}6134 & 0{,}7210 & \\
\bottomrule
\end{tabular}
\label{tab:dak_coef}
\end{table}

A equação~\eqref{eq:DAK} é resolvida iterativamente (ponto fixo):
$Z^{(k+1)} = $ membro direito de~\eqref{eq:DAK} calculado com $Z^{(k)}$,
partindo de $Z^{(0)}=1$, até $|Z^{(k+1)}-Z^{(k)}|<10^{-6}$.

\subsection{Viscosidade do Gás Natural}

\subsubsection{Correlação de Lee, González e Eakin (1966)}

Lee, González e Eakin~\cite{LGE1966} propuseram a seguinte correlação empírica:
\begin{align}
  \mu_g &= K \exp\!\left(X\,\rho_g^Y\right) \times 10^{-4} \quad [\text{cP}]\\
  K &= \frac{(9{,}4 + 0{,}02\,M_g)\,T^{1{,}5}}{209 + 19\,M_g + T}\\
  X &= 3{,}5 + \frac{986}{T} + 0{,}01\,M_g, \qquad Y = 2{,}4 - 0{,}2\,X
\end{align}

onde $M_g = 28{,}97\,\gamma_g$ é o peso molecular do gás e
$\rho_g = M_g P/(10{,}73\,Z\,T)\times 0{,}016018$ é a densidade em
g/cm\textsuperscript{3}.

\subsubsection{Correlação de Lucas et al.\ (1981)}

O método de Lucas et al.~\cite{Lucas1981} calcula a viscosidade a baixa pressão e aplica uma correcção
de pressão~\cite{Ahmed2010}:
\begin{equation}
  \xi = 0{,}176\left(\frac{T_{pc,\text{K}}}{M_g^3\,P_{pc,\text{bar}}^4}\right)^{1/6}
\end{equation}
\begin{equation}
  \mu_1\,\xi = 0{,}807\,T_{pr}^{0{,}618}
               - 0{,}357\,e^{-0{,}449\,T_{pr}}
               + 0{,}340\,e^{-4{,}058\,T_{pr}} + 0{,}018
\end{equation}

A correcção de pressão é aplicada através de quatro parâmetros auxiliares
$a_1$--$a_4$~\cite{Ahmed2010}, obtendo-se a viscosidade final $\mu_g = \mu_1\cdot(\mu_g/\mu_1)$.

\subsection{Factor Volume de Formação e Factor de Expansão}

As expressões seguintes, baseadas na equação de estado real~\cite{Ramos2020}, definem~$B_g$ e~$E_g$:
\begin{align}
  B_g \;[\text{ft}^3/\text{scf}] &= \frac{P_{sc}}{T_{sc}}\cdot\frac{Z\,T}{P}
  \approx 0{,}02827\,\frac{Z\,T}{P}
  \label{eq:Bg}\\[4pt]
  E_g \;[\text{scf/ft}^3] &= \frac{1}{B_g} \approx 35{,}37\,\frac{P}{Z\,T}
  \label{eq:Eg}
\end{align}

\noindent com $P_{sc}=14{,}696\;\text{psia}$ e $T_{sc}=519{,}67\;°\text{R}$
(condições padrão a $60\,°\text{F}$).

% =============================================================================
%  3. METODOLOGIA
% =============================================================================
\section{Metodologia}

\subsection{Arquitectura do Programa}

O programa foi desenvolvido em Python~3.14 com uma arquitectura modular:

\begin{itemize}[noitemsep]
  \item \texttt{modules/properties.py} — correlações de propriedades pseudo-críticas
    e correcções de contaminantes;
  \item \texttt{modules/correlations.py} — factor~$Z$: gás ideal, Hall-Yarborough,
    Dranchuk \& Abou-Kassem;
  \item \texttt{modules/viscosidade.py} — viscosidade: Lee-González-Eakin e Lucas;
  \item \texttt{modules/volumetria.py} — $B_g$ e $E_g$ em diferentes unidades;
  \item \texttt{main.py} — interface gráfica (tkinter/ttk + matplotlib), motor
    de cálculo e integração de todos os módulos.
\end{itemize}

\subsection{Interface Gráfica}

A interface foi desenvolvida com \textit{tkinter} e \textit{ttk}, com tema escuro
moderno. Inclui:
\begin{itemize}[noitemsep]
  \item Painel esquerdo: campos de entrada para pressão, temperatura, densidade,
    contaminantes e selecção de correlações;
  \item Painel direito: tabela de resultados com colunas para $Z_{\text{ideal}}$,
    $Z_{\text{HY}}$, $Z_{\text{DAK}}$, $\mu_g$, $B_g$ e $E_g$;
  \item Janela de gráficos: 4 subgráficos em grelha $2\times2$ — $Z$ vs $P$,
    $\mu_g$ vs $P$, $B_g$ vs $P$ e $E_g$ vs $P$.
\end{itemize}

\subsection{Fluxo de Cálculo}

O cálculo segue os seguintes passos:
\begin{enumerate}[noitemsep]
  \item Leitura e validação das entradas (pressão, temperatura, $\gamma_g$,
    contaminantes, intervalos).
  \item Cálculo de $P_{pc}$ e $T_{pc}$ com a correlação escolhida
    (Standing ou Sutton).
  \item Correcção de $P_{pc}'$ e $T_{pc}'$ com o método escolhido
    (Wichert-Aziz ou Carr-Kobayashi-Burrows).
  \item Para cada valor de pressão na malha:
    \begin{enumerate}[noitemsep, label=\alph*)]
      \item Calcular $P_{pr}$ e $T_{pr}$;
      \item Resolver $Z_{\text{HY}}$ (Newton-Raphson) e
            $Z_{\text{DAK}}$ (ponto fixo);
      \item Calcular $\mu_g$, $B_g$ e $E_g$ com o $Z$ do método seleccionado.
    \end{enumerate}
  \item Exibir resultados em tabela e gráficos.
\end{enumerate}

\subsection{Implementação do Newton-Raphson (Hall-Yarborough)}

O algoritmo de Newton-Raphson foi implementado com as seguintes salvaguardas:
\begin{itemize}[noitemsep]
  \item Valor inicial $y_0 = 0{,}001$ (conforme especificação do docente);
  \item Critério de convergência: $|y^* - y| < 10^{-5}$;
  \item Máximo de 2000 iterações;
  \item Restrição $y > 10^{-7}$ para evitar divisão por zero;
  \item Protecção contra $|f'(y)| < 10^{-15}$.
\end{itemize}

O excerto de código relevante é:
\begin{lstlisting}[caption={Núcleo do método de Newton-Raphson para Hall-Yarborough.},
                   label={lst:newton}]
t = 1.0 / Tpr
A = math.exp(-1.2 * (1.0 - t) ** 2)
y = 0.001                        # condicao inicial

for _ in range(2000):
    a  = -0.06125 * Ppr * t * A
    _b = ((-y + 1.0) * y + 1.0) * y + 1.0
    b  = _b * y / (1.0 - y) ** 3
    c  = -((4.58*t - 9.76)*t + 14.76) * t * y**2
    d  = ((42.4*t - 242.2)*t + 90.7)  * t * y**(2.18 + 2.82*t)
    fy = a + b + c + d

    da  = (((y-4.0)*y+4.0)*y+4.0)*y+1.0
    da /= (1.0 - y) ** 4
    db  = -((9.16*t - 19.52)*t + 29.52) * t * y
    dc  = (2.18 + 2.82*t) * ((42.4*t-242.2)*t+90.7) * t * y**(1.18+2.82*t)
    dfy = da + db + dc

    y_new = y - fy / dfy
    if abs(y_new - y) < 1e-5:
        y = y_new; break
    y = max(y_new, 1e-7)

Z = 0.06125 * Ppr * t * A / y
\end{lstlisting}

% =============================================================================
%  4. RESULTADOS E VALIDAÇÃO
% =============================================================================
\section{Resultados e Validação}

\subsection{Caso de Validação — Gás Doce}

Para um gás doce ($y_{CO_2}=y_{H_2S}=y_{N_2}=0$) com
$\gamma_g=0{,}70$, $T=150\,°\text{F}$ e $P=2000\;\text{psia}$ com
a correlação Standing Seco~\cite{Standing1977}:
\begin{align*}
  P_{pc} &= 669{,}13\;\text{psia}, \quad T_{pc} = 389{,}38\;°\text{R}\\
  T_{pr} &= 1{,}566, \quad P_{pr} = 2{,}989
\end{align*}

\begin{table}[H]
\centering
\caption{Resultados para gás doce, $P=2000\;\text{psia}$, $T=150\;°\text{F}$,
         $\gamma_g=0{,}70$.}
\begin{tabular}{lcc}
\toprule
Propriedade & Valor & Unidade \\
\midrule
$Z_{\text{Ideal}}$   & 1{,}000000 & --- \\
$Z_{\text{HY}}$      & 0{,}810030 & --- \\
$Z_{\text{DAK}}$     & 0{,}811229 & --- \\
$\mu_g$ (LGE)        & 0{,}016921 & cP \\
$B_g$                & 0{,}006983 & ft\textsuperscript{3}/scf \\
$E_g$                & 143{,}21   & scf/ft\textsuperscript{3} \\
\bottomrule
\end{tabular}
\label{tab:val_doce}
\end{table}

\noindent A diferença entre HY e DAK é de $\approx 0{,}15\%$, dentro do
intervalo esperado de 1--2\% para gás doce.

\subsection{Caso de Referência — Gás Ácido (Enunciado)}

Os dados de referência do enunciado correspondem a:
$P=3300\;\text{psia}$, $T=150\;°\text{F}$, $\gamma_g=0{,}75$,
$y_{CO_2}=5\%$, $y_{H_2S}=10\%$, $y_{N_2}=0\%$,
Standing Seco + Wichert-Aziz:

\begin{align*}
  P_{pc}' &= 630{,}07\;\text{psia}, \quad T_{pc}' = 383{,}98\;°\text{R}\\
  T_{pr} &= 1{,}588, \quad P_{pr} = 5{,}238
\end{align*}

\begin{table}[H]
\centering
\caption{Resultados para gás ácido, $P=3300\;\text{psia}$ (caso de referência
         do enunciado).}
\begin{tabular}{lcc}
\toprule
Propriedade & Calculado & Referência \\
\midrule
$Z_{\text{Ideal}}$ & 1{,}000000 & 1{,}0 \\
$Z_{\text{HY}}$    & 0{,}850981 & 0{,}877274\rlap{$^*$} \\
$Z_{\text{DAK}}$   & 0{,}853426 & 0{,}879831\rlap{$^*$} \\
$\mu_g$ (LGE)      & 0{,}023041 & --- \\
$B_g$ (ft\textsuperscript{3}/scf) & 0{,}004446 & --- \\
$E_g$ (scf/ft\textsuperscript{3}) & 224{,}92  & --- \\
\bottomrule
\multicolumn{3}{l}{\small $^*$ Referência obtida com Sutton~(corrigido)+Wichert-Aziz.}\\
\multicolumn{3}{l}{\small \quad Seleccionar ``Gás Natural (Sutton)'' reproduz 0{,}877.}
\end{tabular}
\label{tab:val_acido}
\end{table}

\begin{table}[H]
\centering
\caption{Tabela completa de resultados para $P \in [3000, 3300]\;\text{psia}$
         (passo 50), $T=150\;°\text{F}$, $\gamma_g=0{,}75$, Standing Seco +
         Wichert-Aziz.}
\setlength{\tabcolsep}{6pt}
\begin{tabular}{r r r r r r r}
\toprule
$P$ (psia) & $Z_{\text{ideal}}$ & $Z_{\text{HY}}$ & $Z_{\text{DAK}}$
           & $\mu_g$ (cP) & $B_g$ (ft\textsuperscript{3}/scf) & $E_g$ (scf/ft\textsuperscript{3}) \\
\midrule
3000 & 1{,}000000 & 0{,}833752 & 0{,}836618 & 0{,}021660 & 0{,}004792 & 208{,}70 \\
3050 & 1{,}000000 & 0{,}836322 & 0{,}839141 & 0{,}021892 & 0{,}004728 & 211{,}52 \\
3100 & 1{,}000000 & 0{,}839020 & 0{,}841782 & 0{,}022124 & 0{,}004666 & 214{,}30 \\
3150 & 1{,}000000 & 0{,}841839 & 0{,}844535 & 0{,}022354 & 0{,}004608 & 217{,}03 \\
3200 & 1{,}000000 & 0{,}844776 & 0{,}847396 & 0{,}022584 & 0{,}004552 & 219{,}71 \\
3250 & 1{,}000000 & 0{,}847825 & 0{,}850361 & 0{,}022813 & 0{,}004498 & 222{,}34 \\
3300 & 1{,}000000 & 0{,}850981 & 0{,}853426 & 0{,}023041 & 0{,}004446 & 224{,}92 \\
\bottomrule
\end{tabular}
\label{tab:resultados}
\end{table}

\subsection{Comparação dos Métodos}

A visualização gráfica gerada pela aplicação (janela de gráficos $2\times2$) ilustra o comportamento do factor~$Z$
em função da pressão para os três métodos. Observa-se que:

\begin{itemize}[noitemsep]
  \item O gás ideal ($Z=1$) subestima a desvio real do gás;
  \item Hall-Yarborough e Dranchuk-Abou-Kassem produzem valores muito próximos
    ($\Delta Z < 0{,}3\%$) ao longo de toda a gama de pressões;
  \item $Z$ aumenta ligeiramente com a pressão nesta gama
    ($T_{pr}\approx 1{,}59 > 1$), comportamento típico de gases supercríticos.
\end{itemize}

% =============================================================================
%  5. CONCLUSÕES
% =============================================================================
\section{Conclusões}

O programa desenvolvido satisfaz todos os requisitos do Projecto~Nº~1:

\begin{enumerate}[noitemsep]
  \item Implementação correcta dos três métodos de cálculo do factor~$Z$
    (gás ideal, Hall-Yarborough, Dranchuk \& Abou-Kassem);
  \item Correcções de contaminantes de Wichert-Aziz e Carr-Kobayashi-Burrows;
  \item Correlações de propriedades pseudo-críticas de Standing (seco e húmido)
    e Sutton;
  \item Extensão ao cálculo de viscosidade ($\mu_g$), factor volume de formação
    ($B_g$) e factor de expansão ($E_g$) em diversas unidades;
  \item Interface gráfica moderna com tabela de resultados e quatro gráficos;
  \item Identificação do docente Geraldo Ramos, BSc, MSc, PhD, no cabeçalho;
  \item Executável Windows autónomo gerado por PyInstaller.
\end{enumerate}

A concordância entre Hall-Yarborough e Dranchuk-Abou-Kassem é excelente
($<0{,}3\%$), validando a implementação computacional. O programa constitui
uma ferramenta pedagógica completa para o estudo do comportamento de gases
reais em condições de reservatório.

% =============================================================================
%  REFERÊNCIAS
% =============================================================================
\begin{thebibliography}{9}

\bibitem{HallYarborough1974}
  Hall, K.R.; Yarborough, L. (1974).
  \textit{A new equation of state for Z-factor calculations}.
  Oil \& Gas Journal, 72(25), 82--92.

\bibitem{DAK1975}
  Dranchuk, P.M.; Abou-Kassem, J.H. (1975).
  \textit{Calculation of Z-factors for natural gases using equations of state}.
  Journal of Canadian Petroleum Technology, 14(3), 34--36.

\bibitem{Ramos2020}
  Ramos, G. (2020).
  \textit{Indústria do Gás Natural}, Vol.~2, Cap.~2.
  ISPTEC, Luanda.

\bibitem{Standing1977}
  Standing, M.B. (1977).
  \textit{Volumetric and Phase Behavior of Oil Field Hydrocarbon Systems}.
  Society of Petroleum Engineers of AIME.

\bibitem{Sutton1985}
  Sutton, R.P. (1985).
  \textit{Compressibility Factors for High-Molecular-Weight Reservoir Gases}.
  SPE-14265-MS.

\bibitem{WichertAziz1972}
  Wichert, E.; Aziz, K. (1972).
  \textit{Calculate Z's for sour gases}.
  Hydrocarbon Processing, 51(5), 119--122.

\bibitem{CarrKB1954}
  Carr, N.L.; Kobayashi, R.; Burrows, D.B. (1954).
  \textit{Viscosity of hydrocarbon gas mixtures at high pressure}.
  Transactions of AIME, 201, 264--272.

\bibitem{LGE1966}
  Lee, A.L.; González, M.H.; Eakin, B.E. (1966).
  \textit{The viscosity of natural gas}.
  Journal of Petroleum Technology, 18(8), 997--1000.

\bibitem{Lucas1981}
  Lucas, K. (1981).
  \textit{Die Druckabhängigkeit der Viskosität von Flüssigkeiten — eine einfache Abschätzung}.
  Chemie Ingenieur Technik, 53(12), 959--960.

\bibitem{Ahmed2010}
  Ahmed, T. (2010).
  \textit{Reservoir Engineering Handbook}, 4th~ed.
  Gulf Professional Publishing.

\end{thebibliography}

\end{document}
